\chapter{The Four Fundamental Subspaces}

In Chapter~14, we discovered that every subspace has an unambiguous intrinsic size: its
\textbf{dimension}, counted by the number of vectors in any basis. We saw how Gaussian
elimination systematically partitions the $n$ columns of an $m \times n$ matrix into $r$
pivot columns and $n - r$ free columns. That gave us explicit bases and exact dimension
theorems for two geometric spaces: the column space $\col(\mat{A}) \subseteq \R^m$ with
dimension $r$, and the nullspace $\nullsp(\mat{A}) \subseteq \R^n$ with dimension $n - r$.

Now we look at the matrix through its dual structure: every matrix has rows
as well as columns. Transposing $\mat{A}$ into an $n \times m$ matrix $\mat{A}\T$
introduces two companion subspaces. Together, these form the \emph{four fundamental
subspaces}: two in the input space $\R^n$, and two in the output space $\R^m$.

In each space, the two subspaces are orthogonal complements whose dimensions add
up to the dimension of the ambient space ($r + (n-r) = n$ in $\R^n$, and
$r + (m-r) = m$ in $\R^m$). Gilbert Strang assembled these four spaces into a
unified diagram often called ``the big picture of linear algebra.'' Every matrix
takes an $r$-dimensional subspace of the domain and maps it invertibly onto an
$r$-dimensional subspace of the codomain, while sending the orthogonal complement
to zero.

That is the central geometric structure. The rest of this chapter develops each
space and its properties step by step.

\section{Two spaces for the inputs, two for the outputs}

Let $\mat{A}$ be an $m \times n$ matrix, representing a linear map from $\R^n$ to
$\R^m$. There are two ambient spaces: the input space (domain) $\R^n$ and the
output space (codomain) $\R^m$. Each contains exactly two of the fundamental
subspaces, and in each space the two are \emph{orthogonal complements} of each
other: every vector in one meets every vector in the other at a right angle
(orthogonality), and their dimensions sum to the dimension of the ambient space
(complements).

\begin{focusbox}
The four fundamental subspaces of an $m \times n$ matrix $\mat{A}$ are:
\begin{itemize}
  \item the \textbf{column space} $\col(\mat{A}) \subseteq \R^m$ -- all reachable outputs $\mat{A}\vec{x}$;
  \item the \textbf{nullspace} $\nullsp(\mat{A}) \subseteq \R^n$ -- all inputs mapped to zero by $\mat{A}$;
  \item the \textbf{row space} $\row(\mat{A}) = \col(\mat{A}\T) \subseteq \R^n$ -- the span of the rows;
  \item the \textbf{left nullspace} $\nullsp(\mat{A}\T) \subseteq \R^m$ -- all output directions orthogonal to $\col(\mat{A})$, satisfying $\mat{A}\T\vec{y}=\vec{0}$.
\end{itemize}
Two live in $\R^n$ (column-vector inputs): the row space and the nullspace.
Two live in $\R^m$ (column-vector outputs): the column space and the left nullspace.
\end{focusbox}

\begin{intuitionbox}
A matrix $\mat{A}$ defines a linear mapping from $\R^n$ to $\R^m$. Inside $\R^n$,
every input vector decomposes uniquely into two orthogonal components: one in
the row space and one in the nullspace. $\mat{A}$ maps vectors in the nullspace to
$\vec{0}$, while mapping the row space one-to-one and onto the column space in
$\R^m$. Any vector in $\R^m$ outside the column space lies in the left nullspace
and cannot be reached by $\mat{A}$.
\end{intuitionbox}

Before drawing the big picture, we build each space carefully and learn to read its dimension and
a basis straight off the elimination of $\mat{A}$.

\section{The column space: where the outputs live}

Recall the definition. The column space is everything you can build out of the columns of
$\mat{A}$:
\[
  \col(\mat{A}) \;=\; \set{\,\mat{A}\vec{x} : \vec{x}\in\R^n\,}
  \;=\; \spn(\text{columns of }\mat{A}) \;\subseteq\; \R^m.
\]
The equation $\mat{A}\vec{x}=\vec{b}$ has a solution exactly when $\vec{b}$ lies in this space.
The column space is the set of \emph{achievable right-hand sides}.

\begin{example}\label{ch15:colspace-rank2}
Take
\[
  \mat{A} \;=\;
  \begin{bmatrix} 1 & 2 & 3\\ 2 & 4 & 8\\ 1 & 2 & 5 \end{bmatrix}.
\]
Column~2 is exactly $2\times$ column~1, so it adds nothing new. Columns~1 and~3 point in
different directions. The column space is therefore the plane through the origin in $\R^3$ spanned by
$\colvec{1\\2\\1}$ and $\colvec{3\\8\\5}$. It is two-dimensional: a flat sheet inside three-space.
The vector $\colvec{0\\0\\1}$, say, is \emph{not} on that sheet, so $\mat{A}\vec{x}=\colvec{0\\0\\1}$
has no solution.
\end{example}

How did we know columns~1 and~3 were the ``real'' ones? Elimination tells us. The pivot columns
of $\mat{A}$ -- the columns that end up holding a pivot after row reduction -- are a basis for the
column space.

\begin{formalbox}
\textbf{Basis for the column space.} Row-reduce $\mat{A}$ to find which columns contain pivots.
The corresponding columns \emph{of the original $\mat{A}$} form a basis for $\col(\mat{A})$.
Therefore $\dimn \col(\mat{A}) = r$, the number of pivots, which we call the \textbf{rank}.
\end{formalbox}

\begin{trapbox}
A subtle but classic mistake: after row reduction you must take the pivot columns of the
\emph{original} matrix, not the pivot columns of the reduced matrix. Row operations change the
column space! (They mix the rows, which moves the columns around in $\R^m$.) What row operations
\emph{do} preserve is \emph{which} columns are dependent on which -- so the \emph{positions} of the
pivots are trustworthy, but the actual basis vectors must come from $\mat{A}$ itself.
\end{trapbox}

Here is the trap in miniature.

\begin{labbox}
Reduce $\mat{A}=\begin{bmatrix}1&1\\1&1\end{bmatrix}$. You get
$\begin{bmatrix}1&1\\0&0\end{bmatrix}$. The pivot is in column~1. The column space of the reduced
matrix is the span of $\colvec{1\\0}$ (the $x$-axis). But the column space of $\mat{A}$ is the
span of $\colvec{1\\1}$ (the diagonal line). Different lines! Same dimension ($r=1$), same pivot
positions, but you must read the basis vector $\colvec{1\\1}$ off the original $\mat{A}$.
\end{labbox}

\begin{figure}[H]
\centering
\begin{tikzpicture}[scale=1.1,>=Stealth]
  % axes
  \draw[->,spacegray] (-0.3,0)--(3.2,0) node[right]{$x$};
  \draw[->,spacegray] (0,-0.3)--(0,3.2) node[above]{$y$};
  % column space of A: the diagonal
  \draw[vecblue,very thick] (-0.2,-0.2)--(3,3);
  \fill[vecblue] (1,1) circle (1.6pt);
  \draw[->,vecblue,very thick] (0,0)--(1,1);
  \node[vecblue,right] at (1.05,0.85){$\col(\mat{A})$};
  % column space of reduced: x-axis (dashed red)
  \draw[vecred,very thick,dashed] (-0.2,0)--(3,0);
  \draw[->,vecred,very thick] (0,0)--(1,0);
  \node[vecred,below] at (2.2,-0.05){$\col(\text{reduced})$};
\end{tikzpicture}
\caption{Row reduction preserves pivot \emph{positions} but not the column space itself. For
$\mat{A}=\protect\begin{bmatrix}1&1\\1&1\protect\end{bmatrix}$ the true column space is the
diagonal line (blue); the reduced matrix has the $x$-axis (red dashed). Always read the basis
off the original columns.}
\label{ch15:colspace-trap}
\end{figure}

\begin{recoverybox}
\textbf{Pause and think:} Suppose $\mat{A}$ is a $4 \times 3$ matrix. What is the largest possible
dimension its column space can have? Could it ever fill the output room $\R^4$?
\vspace{0.3em}

The column space is spanned by $3$ columns living in $\R^4$. By Chapter~14's dimension theorem,
$3$ vectors cannot span a $4$-dimensional space: $\dim \col(\mat{A}) \le 3$. So $\col(\mat{A})$ is
at most a 3D subspace inside $\R^4$. There is \emph{always} at least one dimension of unreachable
space left over in the output room.
\end{recoverybox}

\section{The nullspace: inputs mapped to zero}

The nullspace collects every input that $\mat{A}$ sends to zero:
\[
  \nullsp(\mat{A}) \;=\; \set{\,\vec{x}\in\R^n : \mat{A}\vec{x}=\vec{0}\,} \;\subseteq\; \R^n.
\]
This is a subspace because if $\mat{A}\vec{x}=\vec{0}$ and $\mat{A}\vec{y}=\vec{0}$ then
$\mat{A}(c\vec{x}+d\vec{y})=\vec{0}$ too. Geometrically, moving along a nullspace
direction does not change the output: $\mat{A}(\vec{x} + \vec{n}) = \mat{A}\vec{x}$
for any $\vec{n} \in \nullsp(\mat{A})$.

We find a basis by solving $\mat{A}\vec{x}=\vec{0}$ with Gaussian elimination.
Setting each free variable to $1$ while setting all other free variables to $0$
produces one \emph{special solution}, and the special solutions form a basis for
$\nullsp(\mat{A})$.

\begin{example}\label{ch15:nullspace-special}
Continue with $\mat{A}=\begin{bmatrix}1&2&3\\2&4&8\\1&2&5\end{bmatrix}$ from
Example~\ref{ch15:colspace-rank2}. Eliminate:
\[
  \begin{bmatrix}1&2&3\\2&4&8\\1&2&5\end{bmatrix}
  \;\xrightarrow{R_2-2R_1,\,R_3-R_1}\;
  \begin{bmatrix}1&2&3\\0&0&2\\0&0&2\end{bmatrix}
  \;\xrightarrow{R_3-R_2}\;
  \begin{bmatrix}1&2&3\\0&0&2\\0&0&0\end{bmatrix}.
\]
Pivots sit in columns~1 and~3; column~2 is free. Set the free variable $x_2 = 1$. Back-substitute:
the second row gives $2x_3=0$, so $x_3=0$; the first row gives $x_1+2x_2+3x_3=0$, so $x_1=-2$. The
special solution is
\[
  \vec{s} = \colvec{-2\\1\\0},
  \qquad
  \nullsp(\mat{A}) = \spn\set[\big]{\colvec{-2\\1\\0}}.
\]
Check: $\mat{A}\vec{s}=\colvec{1\cdot(-2)+2\cdot1+3\cdot0\\ 2\cdot(-2)+4\cdot1+8\cdot0\\
1\cdot(-2)+2\cdot1+5\cdot0}=\colvec{0\\0\\0}.$ One free variable, so the nullspace is a single
line in $\R^3$: $\dimn\nullsp(\mat{A})=1$.
\end{example}

Notice the bookkeeping that just happened. There were $n=3$ columns. Two held pivots ($r=2$), one
was free ($n-r=1$), and that free variable became the one dimension of the nullspace. This is no
accident.

\begin{formalbox}
\textbf{Dimension of the nullspace.} For an $m\times n$ matrix of rank $r$, the number of free
variables is $n-r$, and each gives an independent special solution. Why are they automatically
independent? Because each special solution has a $1$ in its own free-variable coordinate and a $0$
in all other free-variable coordinates. Chapter~14 guarantees that $n-r$ linearly independent
vectors form a basis of dimension $n-r$. Hence
\[
  \dimn \nullsp(\mat{A}) \;=\; n - r.
\]
This is the \emph{nullity} of $\mat{A}$, and the identity $\rank + \text{nullity} = n$ is the
celebrated \textbf{rank--nullity theorem}.
\end{formalbox}

\begin{walkbox}
Suppose $\mat{A}$ is a $3\times 5$ matrix. On a walk, picture its $5$-dimensional input
space: if $\mat{A}$ crushes a $2$-dimensional plane to zero ($\dimn\nullsp(\mat{A})=2$),
how many independent input directions survive into the output? Could $\nullsp(\mat{A})$
ever be only $1$-dimensional for a $3\times 5$ matrix, or do $3$ rows force at least $2$
free variables?
\end{walkbox}

\section{The row space: inputs that govern the output}

The rows of $\mat{A}$ are vectors with $n$ components, so they live in the domain
$\R^n$. Their span is the row space:
\[
  \row(\mat{A}) \;=\; \spn(\text{rows of }\mat{A}) \;=\; \col(\mat{A}\T) \;\subseteq\; \R^n.
\]
Each component of the output is given by a dot product with a row of $\mat{A}$:
$(\mat{A}\vec{x})_i = (\text{row }i)\cdot \vec{x}$. The output is built entirely
from dot products of $\vec{x}$ against the rows. If $\vec{x}$ is orthogonal to
every row, every dot product is zero and the output vanishes -- that is exactly
the nullspace condition. So the row space and the nullspace are orthogonal
complements inside $\R^n$, a claim we prove in the next section.

How big is the row space, and what is a basis? Here the story is friendlier than for the column
space.

\begin{formalbox}
\textbf{Basis for the row space.} Row operations are reversible combinations of rows, so they
never change the span of the rows. The nonzero rows of the row-reduced matrix $\mat{R}$ span
$\row(\mat{A})$. Why are they linearly independent? Because each nonzero row has a leading pivot
with only zeros below it in that column---the staircase echelon pattern prevents any row from being
canceled by the rows beneath it. Thus the $r$ nonzero rows of $\mat{R}$ form an explicit basis for
$\row(\mat{A})=\row(\mat{R})$, and
\[
  \dimn \row(\mat{A}) \;=\; r \;=\; \dimn\col(\mat{A}).
\]
Unlike the column space, you \emph{may} use the rows of the reduced matrix directly.
\end{formalbox}

That last equation is a striking fact: the number of independent rows equals the number of
independent columns. The next chapter treats it head-on, but you have already
proved it -- both counts equal $r$, the number of pivots.

\begin{example}\label{ch15:rowspace}
Same $\mat{A}$. Its reduced rows (from Example~\ref{ch15:nullspace-special}, after scaling the
pivot to $1$) are
\[
  \row(\mat{A}) = \spn\set[\Big]{\colvec{1\\2\\3},\ \colvec{0\\0\\1}}.
\]
A two-dimensional plane in $\R^3$. And look: the nullspace direction $\colvec{-2\\1\\0}$ is
orthogonal to both row vectors. Dot with $\colvec{1\\2\\3}$: $-2+2+0=0$. Dot with
$\colvec{0\\0\\1}$: $0+0+0=0$. The plane (row space) and the line (nullspace) meet at a perfect
right angle, and together $2+1=3$ they fill all of $\R^3$.
\end{example}

Figure~\ref{ch15:rowspace-null} shows this split of the input room: the row space is a plane, the
nullspace is the line through it at a right angle, and every input vector decomposes uniquely into a
piece in each.

\begin{figure}[H]
\centering
\begin{tikzpicture}[scale=1.0,>=Stealth]
  % a tilted plane (row space) drawn as a parallelogram, with a normal line (nullspace)
  \fill[sourceblue!14,draw=sourceblue!70!black,thick]
        (-2.1,-0.9) -- (1.9,-1.5) -- (2.1,0.9) -- (-1.9,1.5) -- cycle;
  \node[sourceblue!70!black] at (-1.35,1.05) {$\row(\mat{A})$ (plane, $\dim 2$)};
  % nullspace line, perpendicular, drawn rising out of the plane
  \draw[quotientorange!85!black,very thick] (0,0) -- (0.9,2.3);
  \draw[quotientorange!85!black,very thick,dashed] (0,0) -- (-0.9,-2.3);
  \draw[->,quotientorange!85!black,very thick] (0,0) -- (0.7,1.8);
  \node[quotientorange!85!black,right] at (0.7,1.95) {$\nullsp(\mat{A})$ (line, $\dim 1$)};
  % right-angle hint at origin
  \node[fill=black,circle,inner sep=1.1pt] at (0,0) {};
  \node[font=\scriptsize,spacegray,below right] at (0.05,-0.05) {$\vec 0$};
  \node[font=\scriptsize,spacegray] at (0.42,0.4) {$\perp$};
\end{tikzpicture}
\caption{The input room $\R^3$ for the running matrix
(Example~\ref{ch15:rowspace}), split into perpendicular partners: the row space is the plane
spanned by $\protect\colvec{1\\2\\3}$ and $\protect\colvec{0\\0\\1}$, and the nullspace is the line
through $\protect\colvec{-2\\1\\0}$, which stands at a right angle to that plane. Dimensions
$2+1=3$ fill the room.}
\label{ch15:rowspace-null}
\end{figure}

\begin{strangbox}
Strang writes the row space as $C(\mat{A}\T)$ rather than $\row(\mat{A})$ so that all
four subspaces use just two symbols, $C$ and $N$: $C(\mat{A}\T)$ and $N(\mat{A})$ in
$\R^n$, paired with $C(\mat{A})$ and $N(\mat{A}\T)$ in $\R^m$. That notation makes it
immediately visible why transposing $\mat{A}$ swaps the input picture with the output
picture.
\end{strangbox}

\section{The left nullspace: the unreachable output directions}

The fourth space is the one students forget, because naming it requires a transpose. The
\emph{left nullspace} is the nullspace of $\mat{A}\T$:
\[
  \nullsp(\mat{A}\T) \;=\; \set{\,\vec{y}\in\R^m : \mat{A}\T\vec{y}=\vec{0}\,} \;\subseteq\; \R^m.
\]
Why ``left''? Because $\mat{A}\T\vec{y}=\vec{0}$ is the same as $\vec{y}\T\mat{A}=\vec{0}\T$ -- a
row vector $\vec{y}\T$ multiplying $\mat{A}$ from the \emph{left} and producing the zero row. So
$\vec{y}$ records a linear combination of the \emph{rows} of $\mat{A}$ that cancels to zero. These
are exactly the redundancies among the rows.

Geometrically, the left nullspace lives in the output room $\R^m$ and consists of the directions
the column space cannot reach. If $\vec{y}\in\nullsp(\mat{A}\T)$, then for every input $\vec{x}$,
\[
  \vec{y}\cdot(\mat{A}\vec{x}) = \vec{y}\T\mat{A}\vec{x} = (\mat{A}\T\vec{y})\T\vec{x}
  = \vec{0}\T\vec{x}=0.
\]
So $\vec{y}$ is perpendicular to every possible output. The left nullspace is the orthogonal
complement of the column space inside $\R^m$.

\begin{formalbox}
\textbf{Dimension and basis of the left nullspace.} Applying the rank--nullity bookkeeping to
$\mat{A}\T$ (which is $n\times m$ with the same rank $r$):
\[
  \dimn \nullsp(\mat{A}\T) \;=\; m - r.
\]
A clean way to get a basis: do elimination as $\mat{E}\mat{A}=\mat{R}$, tracking the operations
in $\mat{E}$. The rows of $\mat{E}$ that produce the \emph{zero rows} of $\mat{R}$ are exactly a
basis for the left nullspace.
\end{formalbox}

\begin{example}\label{ch15:leftnull}
For our running $\mat{A}$, the reduction $R_2-2R_1$, $R_3-R_1$, then $R_3-R_2$ combined the rows
so that the final third row was all zeros. Tracking the combination, the third row of $\mat{R}$
came from
\[
  (R_3 - R_1) - (R_2 - 2R_1) = R_1 - R_2 + R_3,
\]
i.e. row~1 minus row~2 plus row~3 of $\mat{A}$ gives the zero row. So
$\vec{y}=\colvec{1\\-1\\1}$ satisfies $\vec{y}\T\mat{A}=\vec{0}\T$. Check the first component:
$1\cdot1 - 1\cdot2 + 1\cdot1 = 0$. Second: $1\cdot2-1\cdot4+1\cdot2=0$. Third:
$1\cdot3-1\cdot8+1\cdot5=0$. Indeed $\vec{y}\T\mat{A}=\vec{0}\T$. Here $m=3$, $r=2$, so the left
nullspace is the single line $\spn\set{\colvec{1\\-1\\1}}$ in $\R^3$, dimension $m-r=1$. And it is
orthogonal to the column space plane: dot $\colvec{1\\-1\\1}$ with column~1 $\colvec{1\\2\\1}$ to
get $1-2+1=0$, and with column~3 $\colvec{3\\8\\5}$ to get $3-8+5=0$. Right angles everywhere.
\end{example}

\begin{figure}[H]
\centering
\begin{tikzpicture}[scale=1.0,>=Stealth]
  \fill[targetgreen!14,draw=targetgreen!70!black,thick]
        (-2.1,-0.9) -- (1.9,-1.5) -- (2.1,0.9) -- (-1.9,1.5) -- cycle;
  \node[targetgreen!80!black] at (-1.35,1.05) {$\col(\mat{A})$ (plane, $\dim 2$)};
  \draw[structurepurple!85!black,very thick] (0,0) -- (0.9,2.3);
  \draw[structurepurple!85!black,very thick,dashed] (0,0) -- (-0.9,-2.3);
  \draw[->,structurepurple!85!black,very thick] (0,0) -- (0.7,1.8);
  \node[structurepurple!85!black,right] at (0.7,1.95) {$\nullsp(\mat{A}\T)$ (line, $\dim 1$)};
  \node[fill=black,circle,inner sep=1.1pt] at (0,0) {};
  \node[font=\scriptsize,spacegray,below right] at (0.05,-0.05) {$\vec 0$};
  \node[font=\scriptsize,spacegray] at (0.42,0.4) {$\perp$};
\end{tikzpicture}
\caption{The output room $\R^3$ for the running matrix
(Example~\ref{ch15:leftnull}), split into perpendicular partners: the column space is the plane
spanned by $\protect\colvec{1\\2\\1}$ and $\protect\colvec{3\\8\\5}$, and the left nullspace is the line
through $\protect\colvec{1\\-1\\1}$, standing at a right angle to that plane. Together $2+1=3$
dimensions fill the output room.}
\label{ch15:colspace-leftnull}
\end{figure}

\begin{recoverybox}
\textbf{Pause and think:} If a target output $\vec{b} \in \R^m$ has a nonzero component in the left
nullspace $\nullsp(\mat{A}\T)$, can the linear system $\mat{A}\vec{x} = \vec{b}$ ever have an exact
solution?
\vspace{0.3em}

Never. Every achievable output $\mat{A}\vec{x}$ lives in $\col(\mat{A})$. Because the left nullspace
meets the column space at a right angle, any vector with a piece in $\nullsp(\mat{A}\T)$ sticks out
of the reachable plane. No choice of input $\vec{x}$ can ever reach it.
\end{recoverybox}

\begin{aestheticbox}
The dimensions of the four subspaces exhibit complete symmetry. Because the matrix
and its transpose share the same rank $r$, all four dimensions are determined by just
three integers $m$, $n$, and $r$:
\[
  \dimn\col=r,\quad \dimn\nullsp(\mat{A}\T)=m-r \quad(\text{in }\R^m),\qquad
  \dimn\row=r,\quad \dimn\nullsp=n-r \quad(\text{in }\R^n).
\]
In each room the two dimensions sum to the dimension of the room: $r+(m-r)=m$ and $r+(n-r)=n$.
The whole input space and the whole output space are each split cleanly into a ``used'' part and
an ``unused'' part. A single rank number governs all of it.
\end{aestheticbox}

\section{Orthogonality: the two spaces in each room are perpendicular}

We keep saying ``perpendicular partners.'' Now we prove it, because it is the structural
backbone of the big picture.

\begin{theorem}[Orthogonality of the fundamental subspaces]\label{ch15:ortho-thm}
For any $m\times n$ matrix $\mat{A}$:
\begin{enumerate}
  \item In $\R^n$: every nullspace vector is orthogonal to every row-space vector. In symbols,
  $\nullsp(\mat{A}) \perp \row(\mat{A})$, and together they fill $\R^n$.
  \item In $\R^m$: every left-nullspace vector is orthogonal to every column-space vector. In
  symbols, $\nullsp(\mat{A}\T) \perp \col(\mat{A})$, and together they fill $\R^m$.
\end{enumerate}
\end{theorem}

\begin{proof}
\emph{Part 1.} Let $\vec{x}\in\nullsp(\mat{A})$, so $\mat{A}\vec{x}=\vec{0}$. Reading $\mat{A}\vec{x}$
row by row, the $i$-th entry is $(\text{row}_i)\cdot\vec{x}=0$. So $\vec{x}$ is orthogonal to every
row, hence to every linear combination of rows, hence to the entire row space. That gives
$\nullsp(\mat{A})\perp\row(\mat{A})$.

For the dimension count: $\dimn\row=r$ and $\dimn\nullsp=n-r$. If a vector $\vec{v}$ belongs to
both $\row(\mat{A})$ and $\nullsp(\mat{A})$, it must be orthogonal to itself:
$\vec{v}\cdot\vec{v}=\norm{\vec{v}}^2=0$, forcing $\vec{v}=\vec{0}$. The two subspaces meet only
at the origin. By Chapter~14's direct sum theorem, two subspaces that intersect only at $\vec{0}$
and whose dimensions sum to $r+(n-r)=n=\dimn\R^n$ span the entire room. They are true
orthogonal complements: every vector in $\R^n$ decomposes uniquely into a row-space piece plus a
nullspace piece.

\emph{Part 2} is Part 1 applied to $\mat{A}\T$: its row space is $\col(\mat{A})$ and its nullspace
is $\nullsp(\mat{A}\T)$, both living in $\R^m$.
\end{proof}

\begin{intuitionbox}
Here is the cleanest way to remember which pairs are perpendicular. The nullspace is defined by
``dot with every row gives zero,'' so the nullspace is whatever is left after you remove the row
space -- they are complements. Likewise the left nullspace is ``dot with every column gives zero,''
so it is the complement of the column space. Rows pair with the nullspace; columns pair with the
left nullspace. The output ($\col$) and the unreachable directions ($\nullsp(\mat{A}\T)$) carve up
$\R^m$; the listened-to inputs ($\row$) and the ignored inputs ($\nullsp$) carve up $\R^n$.
\end{intuitionbox}

\begin{connectionbox}
The phrase ``orthogonal complement and the spaces fill the room'' is the
\emph{Fundamental Theorem of Linear Algebra, Part 2}, and it is the finite-dimensional shadow of a
deep theme in analysis. In functional analysis---the branch of mathematics that studies calculus on
infinite-dimensional function spaces---the same statement is the \emph{Fredholm alternative}: the
equation $\mat{A}\vec{x}=\vec{b}$ is solvable precisely when $\vec{b}$ is orthogonal to every
solution of the adjoint homogeneous problem $\mat{A}\T\vec{y}=\vec{0}$ (where ``adjoint'' is the
infinite-dimensional counterpart of the matrix transpose). That adjoint condition is exactly
``$\vec{b}\perp$ left nullspace, i.e. $\vec{b}\in\col(\mat{A})$.'' The four-subspace picture you
are learning here scales all the way up to differential operators and underwrites the solvability
theory for huge classes of differential equations.
\end{connectionbox}

\section{The big picture}

Now we assemble everything into the one diagram worth memorizing. Draw the input room $\R^n$ on
the left and the output room $\R^m$ on the right. Inside each room, draw the two perpendicular
subspaces, and show the action of $\mat{A}$ as arrows crossing from left to right.

\begin{figure}[H]
\centering
\begin{tikzpicture}[scale=0.96,>=Stealth]
  % ===== DOMAIN R^n (LEFT) =====
  \draw[fill=gray!2, draw=gray!25, rounded corners=4pt] (-6.4,-3.2) rectangle (-1.2,3.8);
  \node[draw=spacegray!60, fill=gray!10, rounded corners=2pt, font=\footnotesize\scshape, text=black!85]
    at (-3.8, 3.4) {Domain $\R^n$ (Inputs)};

  % Row space card (top)
  \draw[fill=sourceblue!12, draw=sourceblue!70, rounded corners=3pt] (-6.1,0.25) rectangle (-1.5,2.95);
  \node[draw=sourceblue!80, fill=sourceblue!20, rounded corners=2pt, font=\scriptsize\bfseries, text=sourceblue!95!black]
    at (-3.8, 2.6) {Row Space $\row(\mat{A}) = \col(\mat{A}^T)$};
  \node[font=\scriptsize, text=sourceblue!85!black] at (-3.8, 2.1) {$\dim = r$};
  \draw[->, very thick, vecblue] (-3.8,0.25) -- (-2.4,1.4) node[above left, font=\scriptsize] {$\vec{x}_{\mathrm{row}}$};
  \fill[vecblue] (-2.4,1.4) circle (1.6pt);

  % Nullspace card (bottom)
  \draw[fill=quotientorange!12, draw=quotientorange!70, rounded corners=3pt] (-6.1,-2.95) rectangle (-1.5,-0.25);
  \node[draw=quotientorange!80, fill=quotientorange!20, rounded corners=2pt, font=\scriptsize\bfseries, text=quotientorange!95!black]
    at (-3.8, -0.6) {Nullspace $\nullsp(\mat{A})$};
  \node[font=\scriptsize, text=quotientorange!85!black] at (-3.8, -1.05) {$\dim = n - r$};
  \draw[->, very thick, quotientorange!90!black] (-3.8,-0.25) -- (-5.2,-1.5) node[below left, font=\scriptsize] {$\vec{x}_{\mathrm{null}}$};
  \fill[quotientorange!90!black] (-5.2,-1.5) circle (1.6pt);

  % Domain Origin & Orthogonality
  \fill[black] (-1.65,0) circle (2pt) node[left=1pt, font=\scriptsize] {$\vec{0}$};
  \node[font=\tiny, spacegray] at (-4.2, 0) {$\row \perp \nullsp$};

  % ===== CODOMAIN R^m (RIGHT) =====
  \draw[fill=gray!2, draw=gray!25, rounded corners=4pt] (1.2,-3.2) rectangle (6.4,3.8);
  \node[draw=spacegray!60, fill=gray!10, rounded corners=2pt, font=\footnotesize\scshape, text=black!85]
    at (3.8, 3.4) {Codomain $\R^m$ (Outputs)};

  % Column space card (top)
  \draw[fill=targetgreen!12, draw=targetgreen!70, rounded corners=3pt] (1.5,0.25) rectangle (6.1,2.95);
  \node[draw=targetgreen!80, fill=targetgreen!20, rounded corners=2pt, font=\scriptsize\bfseries, text=targetgreen!95!black]
    at (3.8, 2.6) {Column Space $\col(\mat{A})$};
  \node[font=\scriptsize, text=targetgreen!85!black] at (3.8, 2.1) {$\dim = r$};
  \draw[->, very thick, targetgreen!80!black] (3.8,0.25) -- (5.2,1.4) node[above right, font=\scriptsize] {$\vec{b} = \mat{A}\vec{x}_{\mathrm{row}}$};
  \fill[targetgreen!80!black] (5.2,1.4) circle (1.6pt);

  % Left nullspace card (bottom)
  \draw[fill=structurepurple!12, draw=structurepurple!70, rounded corners=3pt] (1.5,-2.95) rectangle (6.1,-0.25);
  \node[draw=structurepurple!80, fill=structurepurple!20, rounded corners=2pt, font=\scriptsize\bfseries, text=structurepurple!95!black]
    at (3.8, -0.6) {Left Nullspace $\nullsp(\mat{A}^T)$};
  \node[font=\scriptsize, text=structurepurple!85!black] at (3.8, -1.05) {$\dim = m - r$};
  \node[font=\tiny, spacegray, align=center] at (3.8, -1.8) {Orthogonal to all outputs\\[2pt] $\mat{A}^T\vec{y} = \vec{0}$};

  % Codomain Origin & Orthogonality
  \fill[black] (1.65,0) circle (2pt) node[right=1pt, font=\scriptsize] {$\vec{0}$};
  \node[font=\tiny, spacegray] at (4.2, 0) {$\col \perp \nullsp(\mat{A}^T)$};

  % ===== MAPPINGS ACROSS SPACES =====
  % Main isomorphism: row space -> column space
  \draw[->, line width=1.6pt, sourceblue] (-1.5, 1.55) -- (1.5, 1.55)
    node[midway, above, font=\small\bfseries] {$\mat{A}$}
    node[midway, below, font=\tiny, align=center] {invertible bijection\\[1pt]($r$-dim onto $r$-dim)};

  % Reverse mapping A^T: column space -> row space
  \draw[->, thick, dashed, targetgreen!80!black] (1.5, 0.55) -- (-1.5, 0.55)
    node[midway, below, font=\tiny] {$\mat{A}^T : \col(\mat{A}) \to \row(\mat{A})$};

  % Nullspace squash: nullspace -> codomain origin (stays in center corridor)
  \draw[->, line width=1.3pt, quotientorange!85!black] (-1.5, -1.05) .. controls (0.1, -1.05) and (0.6, -0.05) .. (1.52, -0.02);
  \node[quotientorange!95!black, font=\tiny, fill=white, inner sep=1pt] at (0.0, -1.25) {$\mat{A}(\nullsp) = \{\vec{0}\}$};

  % Left nullspace squash: left nullspace -> domain origin (stays in center corridor)
  \draw[->, thick, dashed, structurepurple!85!black] (1.5, -2.05) .. controls (-0.1, -2.05) and (-0.6, -0.05) .. (-1.52, -0.02);
  \node[structurepurple!95!black, font=\tiny, fill=white, inner sep=1pt] at (0.0, -2.35) {$\mat{A}^T(\nullsp(\mat{A}^T)) = \{\vec{0}\}$};
\end{tikzpicture}
\caption{The Big Picture of Linear Algebra. The input domain $\R^n$ decomposes into the row space
and nullspace (orthogonal complements of dimensions $r$ and $n-r$). The codomain $\R^m$ decomposes
into the column space and left nullspace (orthogonal complements of dimensions $r$ and $m-r$).
The matrix $\mat{A}$ maps the entire nullspace directly to the origin $\vec{0} \in \R^m$,
while establishing an invertible, one-to-one isomorphism between the $r$-dimensional row space
and the $r$-dimensional column space.}
\label{ch15:bigpicture}
\end{figure}

Two facts drive everything in Figure~\ref{ch15:bigpicture}. First, \textbf{the nullspace maps to zero}
by definition. Second, \textbf{the row space maps invertibly onto the column space}. Take any input $\vec{x}\in\R^n$ and split it (uniquely, by
Theorem~\ref{ch15:ortho-thm}) as
\[
  \vec{x} = \vec{x}_{\text{row}} + \vec{x}_{\text{null}},
  \qquad \vec{x}_{\text{row}}\in\row(\mat{A}),\ \ \vec{x}_{\text{null}}\in\nullsp(\mat{A}).
\]
Then
\[
  \mat{A}\vec{x} = \mat{A}\vec{x}_{\text{row}} + \mat{A}\vec{x}_{\text{null}}
  = \mat{A}\vec{x}_{\text{row}} + \vec{0} = \mat{A}\vec{x}_{\text{row}}.
\]
So the output is entirely determined by the row-space component of the input: the nullspace
component does not affect the output. The map ``$\vec{x}_{\text{row}}\mapsto\mat{A}\vec{x}_{\text{row}}$''
goes from the row space (dimension $r$) to the column space (dimension $r$).

An \textbf{isomorphism} is a linear map that is both one-to-one (different inputs never collide)
and onto (every target vector in the space is reached)---an invertible linear map between two
vector spaces that preserves vector addition and scalar multiplication.

\begin{theorem}[Row space to column space is an isomorphism]\label{ch15:iso-thm}
Restricted to the row space, $\mat{A}$ is an isomorphism from $\row(\mat{A})$ onto
$\col(\mat{A})$: every output in $\col(\mat{A})$ comes from exactly one vector in $\row(\mat{A})$.
\end{theorem}

\begin{proof}
\emph{Onto:} every output $\mat{A}\vec{x}$ equals $\mat{A}\vec{x}_{\text{row}}$ by the splitting
above, so every element of the column space is hit by some row-space vector.

\emph{One-to-one:} suppose two row-space vectors $\vec{u},\vec{v}\in\row(\mat{A})$ give the same
output, $\mat{A}\vec{u}=\mat{A}\vec{v}$. Then $\mat{A}(\vec{u}-\vec{v})=\vec{0}$, so
$\vec{u}-\vec{v}\in\nullsp(\mat{A})$. But $\vec{u}-\vec{v}$ is also in the row space (a difference
of row-space vectors). A vector in both the row space and the nullspace is orthogonal to itself,
so $\vec{u}-\vec{v}=\vec{0}$, i.e. $\vec{u}=\vec{v}$. Distinct row-space vectors give distinct
outputs.
\end{proof}

\begin{focusbox}
Restricting the domain to the row space and the codomain to the column space makes every linear
map invertible. By removing the nullspace from the input space and the left nullspace from the
output space, $\mat{A}$ defines a bijective linear transformation from $\row(\mat{A})$ to
$\col(\mat{A})$. The rank $r$ is the intrinsic dimension of this invertible mapping.
\end{focusbox}

\begin{example}\label{ch15:splitting}
Watch the splitting in action with our $\mat{A}$. Take the input $\vec{x}=\colvec{0\\1\\0}$
(just turning on the free variable, plus nothing else). Decompose it. The nullspace direction is
$\vec{n}=\colvec{-2\\1\\0}$ with $\norm{\vec{n}}^2=5$. Project $\vec{x}$ onto it:
\[
  \vec{x}_{\text{null}} = \frac{\vec{x}\cdot\vec{n}}{\vec{n}\cdot\vec{n}}\,\vec{n}
  = \frac{(0)(-2)+(1)(1)+0}{5}\colvec{-2\\1\\0}
  = \frac{1}{5}\colvec{-2\\1\\0}.
\]
Then $\vec{x}_{\text{row}} = \vec{x}-\vec{x}_{\text{null}}
= \colvec{0\\1\\0}-\tfrac{1}{5}\colvec{-2\\1\\0} = \tfrac{1}{5}\colvec{2\\4\\0}$. You can verify
$\vec{x}_{\text{row}}$ lies in the row space (it is $\tfrac{2}{5}\colvec{1\\2\\0}$, a combination of
the rows) and is orthogonal to $\vec{n}$. Now apply $\mat{A}$:
\[
  \mat{A}\vec{x} = \mat{A}\colvec{0\\1\\0}=\colvec{2\\4\\2},\qquad
  \mat{A}\vec{x}_{\text{row}} = \mat{A}\cdot\tfrac15\colvec{2\\4\\0}
  =\tfrac15\colvec{2\cdot1+4\cdot2\\2\cdot2+4\cdot4\\2\cdot1+4\cdot2}=\tfrac15\colvec{10\\20\\10}
  =\colvec{2\\4\\2}.
\]
Identical. The nullspace part contributed nothing; the row-space part did all the work. That is
Theorem~\ref{ch15:iso-thm} with numbers.
\end{example}

Figure~\ref{ch15:split-arrows} illustrates this decomposition: the input $\vec{x}$ projects
orthogonally onto the row space and onto the nullspace. $\mat{A}$ maps the nullspace
component to zero, so the output $\mat{A}\vec{x}$ depends entirely on $\vec{x}_{\text{row}}$.

\begin{figure}[H]
\centering
\begin{tikzpicture}[scale=1.0,>=Stealth,font=\small]
  % input room
  \node[spacegray] at (-2.0,2.5) {input room};
  \draw[sourceblue!70!black,very thick] (-3.6,-1.4) -- (-0.4,1.4);
  \node[sourceblue!70!black,right] at (-1.0,1.25) {row space};
  \draw[quotientorange!85!black,very thick,dashed] (-3.4,1.2) -- (-0.6,-1.2);
  \node[quotientorange!85!black,left] at (-3.35,1.15) {nullspace};
  \coordinate (O) at (-2.0,0);
  \node[fill=black,circle,inner sep=1pt] at (O) {};
  % x and its two pieces
  \coordinate (xrow) at (-1.0,0.75);   % along row space dir
  \coordinate (xnul) at (-2.55,0.75);  % along nullspace dir
  \coordinate (X) at ($(xrow)+(xnul)-(O)$);
  \draw[->,black,very thick] (O) -- (X) node[above right]{$\vec{x}$};
  \draw[->,sourceblue!70!black,thick] (O) -- (xrow) node[below right=-1pt]{$\vec{x}_{\text{row}}$};
  \draw[->,quotientorange!85!black,thick] (O) -- (xnul) node[above left=-2pt]{$\vec{x}_{\text{null}}$};
  \draw[gray!60,dashed] (X)--(xrow); \draw[gray!60,dashed] (X)--(xnul);
  % arrow A
  \draw[->,structurepurple,very thick] (0.6,0.3) -- (2.0,0.3) node[midway,above]{$\mat{A}$};
  % output room
  \node[spacegray] at (4.2,2.5) {output room};
  \draw[targetgreen!80!black,very thick] (2.6,-1.2) -- (5.6,1.4);
  \node[targetgreen!80!black,right] at (4.2,1.2) {column space};
  \coordinate (O2) at (3.4,-0.2);
  \node[fill=black,circle,inner sep=1pt] at (O2) {};
  \coordinate (AX) at (4.9,1.0);
  \draw[->,targetgreen!80!black,very thick] (O2) -- (AX) node[right]{$\mat{A}\vec{x}=\mat{A}\vec{x}_{\text{row}}$};
\end{tikzpicture}
\caption{The orthogonal decomposition of Example~\ref{ch15:splitting}. The input $\vec{x}$ resolves
into a row-space component $\vec{x}_{\text{row}}$ and a nullspace component $\vec{x}_{\text{null}}$.
Because $\mat{A}\vec{x}_{\text{null}} = \vec{0}$, the output depends entirely on $\vec{x}_{\text{row}}$
and lies in the column space.}
\label{ch15:split-arrows}
\end{figure}

\section{Reading all four bases off one elimination}

One row reduction hands you everything, so it is worth consolidating the recipe in one place.

\begin{formalbox}
\textbf{The one-elimination recipe.} Reduce $\mat{A}$ (tracking operations in $\mat{E}$ so that
$\mat{E}\mat{A}=\mat{R}$). Let the pivots sit in $r$ columns.
\begin{itemize}
  \item \textbf{Column space} ($\dim r$, in $\R^m$): the pivot columns of the \emph{original}
  $\mat{A}$.
  \item \textbf{Row space} ($\dim r$, in $\R^n$): the nonzero rows of $\mat{R}$.
  \item \textbf{Nullspace} ($\dim n-r$, in $\R^n$): the special solutions, one per free column.
  \item \textbf{Left nullspace} ($\dim m-r$, in $\R^m$): the rows of $\mat{E}$ that produced the
  zero rows of $\mat{R}$.
\end{itemize}
\end{formalbox}

\begin{labbox}
Try the full recipe on
\[
  \mat{A} = \begin{bmatrix} 1 & 2 & 1 & 1\\ 2 & 4 & 3 & 4\\ 1 & 2 & 2 & 3 \end{bmatrix}
  \quad (m=3,\ n=4).
\]
Eliminate to $\mat{R}=\begin{bmatrix}1&2&0&-1\\0&0&1&2\\0&0&0&0\end{bmatrix}$ (pivots in columns 1
and 3, so $r=2$). Now read off all four:
\begin{itemize}
  \item \textbf{Column space}: columns 1 and 3 of $\mat{A}$, namely $\colvec{1\\2\\1}$ and
  $\colvec{1\\3\\2}$. ($\dim 2$, in $\R^3$.)
  \item \textbf{Row space}: $\colvec{1\\2\\0\\-1}$ and $\colvec{0\\0\\1\\2}$. ($\dim 2$, in
  $\R^4$.)
  \item \textbf{Nullspace}: free columns are 2 and 4. Setting $x_2=1$ (rest free $=0$) gives
  $\colvec{-2\\1\\0\\0}$; setting $x_4=1$ gives $\colvec{1\\0\\-2\\1}$. ($\dim 2$, in $\R^4$.)
  \item \textbf{Left nullspace}: one zero row, so $\dim 1$ in $\R^3$. The combination
  $R_1-R_2+R_3$ of the original rows is zero, giving $\colvec{1\\-1\\1}$.
\end{itemize}
Sanity check the dimensions: $r=2$, $n-r=2$ (fills $\R^4$), $m-r=1$ (with $r=2$ fills $\R^3$).
\end{labbox}

\begin{trapbox}
Watch your rooms. The row space and nullspace are sets of $4$-component vectors (they live in the
\emph{input} space $\R^4$); the column space and left nullspace are $3$-component vectors (they live
in the \emph{output} space $\R^3$). A frequent error is to try to ``check orthogonality'' between
the column space and the nullspace -- but those two live in different rooms and cannot even be
dotted together. Orthogonality is always \emph{within} a room: row $\perp$ null in $\R^n$, column
$\perp$ left-null in $\R^m$.
\end{trapbox}

\section{A modern application: where data lives and what is unreachable}

\begin{applicationbox}
\textbf{Recommender systems and the column space.} Imagine a big matrix $\mat{A}$ whose rows are
movies and whose columns are users, with entry $A_{ij}$ the rating user $j$ gave movie $i$. Most
entries are missing; the engine's job is to fill them in. The working hypothesis behind systems
like the ones that won the Netflix Prize is that tastes are \emph{low rank}: a handful of hidden
factors (how much you like comedy, action, slow burns) explain most of the ratings. ``Low rank''
means the rank $r$ is small compared with $m$ and $n$.

The four subspaces tell you exactly what such a model can and cannot do.
\begin{itemize}
  \item The \textbf{column space} is the set of rating patterns the model can produce. If a real
  user's true ratings lie outside this $r$-dimensional column space, no choice of weights can
  reproduce them -- their genuinely individual taste is, mathematically, in the
  \textbf{left nullspace}, perpendicular to everything the model can represent. That orthogonal
  leftover is precisely the part the model must give up on.
  \item The \textbf{nullspace} is the model's blind spot among inputs: directions in
  user/feature space you can change without changing a single predicted rating. Regularizers and
  extra data are how engineers try to shrink or pin down that nullspace so predictions become
  determined rather than free.
\end{itemize}
The same decomposition powers image compression, topic models for text, and the embeddings inside
modern recommendation and language models: keep the $r$ directions that carry signal (the column
space), discard the rest. The rank is the budget; the four subspaces are the accounting.
\end{applicationbox}

\begin{applicationbox}
\textbf{Least squares and the left nullspace, previewed.} When $\mat{A}\vec{x}=\vec{b}$ has no
solution -- the overdetermined case at the heart of every regression -- it is precisely because
$\vec{b}$ has a piece sticking out of the column space, a piece lying in the left nullspace. The
least-squares method, which you will meet soon, throws that unreachable piece away by projecting
$\vec{b}$ onto the column space and solving for the part that \emph{is} achievable. The left
nullspace is the exact measure of ``how much your data refuses to
fit your model.''
\end{applicationbox}

\section{Worked summary example, all four at once}

\begin{example}\label{ch15:full-worked}
Let
\[
  \mat{A} = \begin{bmatrix} 1 & 3 & 2\\ 2 & 6 & 4 \end{bmatrix}
  \qquad (m=2,\ n=3).
\]
Row~2 is twice row~1, so we expect $r=1$. Eliminate: $R_2-2R_1$ gives
$\mat{R}=\begin{bmatrix}1&3&2\\0&0&0\end{bmatrix}$. One pivot, $r=1$.
\begin{itemize}
  \item \textbf{Column space} (in $\R^2$, $\dim 1$): pivot column~1 of $\mat{A}$ is
  $\colvec{1\\2}$. The column space is the line through $\colvec{1\\2}$.
  \item \textbf{Left nullspace} (in $\R^2$, $\dim m-r=1$): the operation $R_2-2R_1$ zeroed a row,
  so the combination $-2(\text{row}_1)+1(\text{row}_2)$ vanishes, i.e. $\vec{y}=\colvec{-2\\1}$
  gives $\vec{y}\T\mat{A}=\vec{0}\T$. Check entry by entry: $-2(1)+1(2)=0$, $-2(3)+1(6)=0$, and
  $-2(2)+1(4)=0$. Good. The left nullspace is the line through $\colvec{-2\\1}$, perpendicular to
  the column-space direction $\colvec{1\\2}$ (since $(-2)(1)+(1)(2)=0$). Together they fill
  $\R^2$.
  \item \textbf{Row space} (in $\R^3$, $\dim 1$): the nonzero reduced row $\colvec{1\\3\\2}$. The
  line through it.
  \item \textbf{Nullspace} (in $\R^3$, $\dim n-r=2$): free columns 2, 3. Special solutions
  $\colvec{-3\\1\\0}$ (set $x_2=1$) and $\colvec{-2\\0\\1}$ (set $x_3=1$). A plane. Both are
  orthogonal to the row space direction $\colvec{1\\3\\2}$: $-3+3+0=0$ and $-2+0+2=0$. The line
  (row space) plus the plane (nullspace) fill $\R^3$.
  \end{itemize}
Every dimension checks: $\R^2 = 1 + 1$, $\R^3 = 1 + 2$. A rank-one matrix: all outputs
are scalar multiples of a single column vector, and all row vectors are multiples of a
single row vector.
\end{example}

\begin{figure}[H]
\centering
\begin{tikzpicture}[scale=1.0,>=Stealth]
  % R^2 output room
  \begin{scope}
    \draw[->,spacegray] (-2.2,0)--(2.2,0) node[right]{$y_1$};
    \draw[->,spacegray] (0,-2.2)--(0,2.2) node[left]{$y_2$};
    \node[spacegray] at (0,2.6) {output room $\R^2$};
    % column space line through (1,2)
    \draw[targetgreen!80!black,very thick] (-1,-2)--(1,2);
    \draw[->,targetgreen!80!black,very thick] (0,0)--(1,2);
    \node[targetgreen!80!black,right] at (0.8,1.7){$\col(\mat{A})$};
    % left null line through (-2,1)
    \draw[structurepurple,very thick,dashed] (2,-1)--(-2,1);
    \draw[->,structurepurple,very thick] (0,0)--(-1.4,0.7);
    \node[structurepurple!85!black,left] at (-1.3,0.95){$\nullsp(\mat{A}\T)$};
    \node[fill=black,circle,inner sep=1.1pt] at (0,0) {};
    \node[font=\scriptsize,spacegray,below right] at (0.05,-0.05) {$\vec{0}$};
    \node[font=\scriptsize,spacegray] at (0.35,0.35) {$\perp$};
  \end{scope}
\end{tikzpicture}
\caption{The output room $\R^2$ for the rank-one matrix of
Example~\ref{ch15:full-worked}: the column space (green, the line through
$\protect\colvec{1\\2}$) and the left nullspace (purple dashed, the perpendicular line through
$\protect\colvec{-2\\1}$). They meet at a right angle and together fill the plane.}
\label{ch15:rank1-room}
\end{figure}

\begin{historybox}
For nearly a century, from Arthur Cayley and James Joseph Sylvester's work in the 1850s
through the mid-twentieth century, linear algebra was taught almost exclusively through the
lens of determinants (numbers computed by cross-multiplying matrix entries to test if an
$n\times n$ matrix is invertible), minors (determinants of smaller square submatrices), and
Cramer's rule. Rectangular $m \times n$ matrices were treated as clumsy anomalies because you
cannot compute the determinant of a non-square matrix. Students spent weeks expanding tedious
polynomial formulas by hand, learning almost nothing about what a transformation actually does
to space.

The geometric breakthrough arrived from an unexpected direction: continuous physics. In 1903,
the Swedish mathematician Ivar Fredholm was studying integral equations---equations where the
unknown is an entire function rather than a list of numbers, describing physical phenomena like
heat distribution and electrostatics. These act like infinite-dimensional linear systems, where
determinants fail completely. Fredholm discovered what is now celebrated as the \emph{Fredholm
Alternative}: a linear system $\mat{A}\vec{x}=\vec{b}$ has an exact solution if and only if
the target $\vec{b}$ is perpendicular to every solution of the transposed equation
$\mat{A}\T\vec{y}=\vec{0}$. In geometric terms, the outputs you can reach ($\col(\mat{A})$)
are the exact orthogonal complement of the vectors killed by the transpose:
\[
  \col(\mat{A}) \;=\; (\nullsp(\mat{A}\T))^\perp.
\]
Whether a system can be solved is a test of right angles.

Despite Fredholm's discovery, textbooks clung to determinant formulas for decades. A major
stumbling block was the transpose itself: students and early calculators were taught to view
$\mat{A}\T$ as a trivial typographical trick---simply ``flipping rows into columns.'' This
blind spot hid the deeper truth: the transpose represents the \emph{adjoint} map, the machine
that translates measurements in the output room back into coordinates in the input room. A second
confusion arose from row reduction: row operations preserve the row space and nullspace in the
input room $\R^n$, but they distort the column space and left nullspace in the output room $\R^m$.

The pedagogical revolution that placed geometry front and center was led by Gilbert Strang. In a
landmark 1993 paper in the \emph{American Mathematical Monthly} titled ``The Fundamental Theorem
of Linear Algebra,'' Strang championed the two-room diagram---two perpendicular subspaces in
$\R^n$, two perpendicular subspaces in $\R^m$, joined by an invertible $r$-dimensional
bridge---as the organizing blueprint of the entire discipline.

This four-subspace picture organizes several central tools of applied mathematics: least-squares data
fitting (projecting an unreachable data vector onto the column space), the Moore--Penrose
pseudoinverse, and the Singular Value Decomposition (SVD), which finds ideal perpendicular bases
for all four subspaces simultaneously. Today, an active frontier lies in high dimensions and massive
datasets: when data matrices have millions of noisy entries, the exact boundary between the row
space and the nullspace blurs. Modern researchers rely on randomized numerical linear algebra
and perturbation bounds (such as the Davis--Kahan theorem) to reliably estimate these four
subspaces. At the same time, extending this clean four-subspace decomposition to higher-order
\emph{tensors} (multi-dimensional grids of data beyond 2D matrices) remains notoriously difficult:
while matrix rank is computed in milliseconds via elimination, computing the rank of a 3D tensor
is provably NP-hard.
\end{historybox}

\section{What we now know}

\begin{focusbox}
The four fundamental subspaces of an $m\times n$ matrix $\mat{A}$ of rank $r$:
\begin{center}
\begin{tabular}{lll}
\toprule
Subspace & Lives in & Dimension \\
\midrule
Column space $\col(\mat{A})$ & $\R^m$ & $r$ \\
Left nullspace $\nullsp(\mat{A}\T)$ & $\R^m$ & $m-r$ \\
Row space $\row(\mat{A})$ & $\R^n$ & $r$ \\
Nullspace $\nullsp(\mat{A})$ & $\R^n$ & $n-r$ \\
\bottomrule
\end{tabular}
\end{center}
In $\R^n$: row space $\perp$ nullspace, dimensions add to $n$.
In $\R^m$: column space $\perp$ left nullspace, dimensions add to $m$.
And $\mat{A}$ maps the row space one-to-one onto the column space, collapsing the nullspace to
$\vec{0}$.
\end{focusbox}

The next chapter gives the rank itself center stage and states the
Fundamental Theorem of Linear Algebra in full -- you have already done most of the work.

\begin{connectionbox}
We now possess the complete structural blueprint for any linear transformation $\mat{A}: \R^n \to \R^m$. Both the input space $\R^n$ and output space $\R^m$ split cleanly into orthogonal complementary subspaces. Every input decomposes into a part that $\mat{A}$ transmits invertibly onto $\col(\mat{A})$ and a part it maps to $\vec{0}$ in $\nullsp(\mat{A})$.

While the Big Picture reveals orthogonality and matching dimensions ($r$ and $n-r$, $r$ and $m-r$), it is still qualitative. It does not provide the explicit metric projection matrices to drop a vector onto $\col(\mat{A})$ or $\row(\mat{A})$. And when $\vec{b} \notin \col(\mat{A})$, knowing that the error lives in $\nullsp(\mat{A}\T)$ does not yet tell us how to compute the best approximate solution. In Chapter~16, we prove that row rank always equals column rank ($\dim \row(\mat{A}) = \dim \col(\mat{A})$) for every matrix without exception, culminating in the complete Fundamental Theorem of Linear Algebra and showing how any matrix decomposes into a sum of rank-one building blocks. That sets the stage for Part~IV (Chapters 17--20), where we construct orthogonal projections and least squares regression.
\end{connectionbox}

\section{Exercises}

\begin{exercise}[Mental warmup: room and dimension check]\label{ch15:ex-mental-dim}
A matrix $\mat{A}$ has size $4 \times 6$ and rank $r = 3$.
\begin{enumerate}[label=(\alph*)]
  \item In your head, state the dimension of each of the four fundamental subspaces.
  \item Which two subspaces live in the input room $\R^6$, and what is the sum of their dimensions?
  \item Which two subspaces live in the output room $\R^4$, and what is the sum of their dimensions?
\end{enumerate}
\end{exercise}

\begin{exercise}[Mental warmup: which room?]\label{ch15:ex-mental-room}
Let $\mat{A}$ be a $3 \times 5$ matrix. A student writes: ``We checked that $\col(\mat{A})$ is
orthogonal to $\nullsp(\mat{A})$ by computing the dot product of their basis vectors.''
Without doing any arithmetic, explain in one sentence why this calculation is impossible.
\end{exercise}

\begin{exercise}\label{ch15:ex-basic}
For $\mat{A}=\begin{bmatrix}1&2&0\\0&0&1\end{bmatrix}$, find a basis and the dimension of each of
the four fundamental subspaces. State which room ($\R^2$ or $\R^3$) each lives in.
\end{exercise}

\begin{exercise}\label{ch15:ex-special}
Let $\mat{A}=\begin{bmatrix}1&1&2&1\\2&2&4&3\\1&1&2&0\end{bmatrix}$. Reduce it, identify the
pivot and free columns, and write down all the special solutions (a basis for the nullspace).
What is the nullity?
\end{exercise}

\begin{exercise}\label{ch15:ex-colspace-trap}
Give a $2\times2$ matrix whose row-reduced form has a pivot in column~1, yet whose column space is
\emph{not} the $x$-axis. Use it to explain in one sentence why the column-space basis must be read
from the original matrix, not the reduced one.
\end{exercise}

\begin{exercise}\label{ch15:ex-leftnull}
For $\mat{A}=\begin{bmatrix}1&2\\2&4\\3&6\end{bmatrix}$, find the left nullspace directly by
solving $\mat{A}\T\vec{y}=\vec{0}$. Confirm its dimension equals $m-r$ and that every basis vector
is orthogonal to the column space.
\end{exercise}

\begin{exercise}\label{ch15:ex-dimcheck}
A matrix $\mat{A}$ is $5\times 7$ with rank $4$. State the dimension of each of the four
fundamental subspaces and the room each lives in. How many independent special solutions does
$\mat{A}\vec{x}=\vec{0}$ have?
\end{exercise}

\begin{exercise}\label{ch15:ex-splitting}
For $\mat{A}=\begin{bmatrix}1&1\\1&1\end{bmatrix}$ and input $\vec{x}=\colvec{3\\1}$, split
$\vec{x}$ into its row-space and nullspace pieces and verify that $\mat{A}\vec{x}=\mat{A}\vec{x}_
{\text{row}}$, the content of Theorem~\ref{ch15:iso-thm}.
\end{exercise}

\begin{exercise}\label{ch15:ex-solvability}
Use the left nullspace to decide solvability. For
$\mat{A}=\begin{bmatrix}1&2\\2&4\\0&1\end{bmatrix}$, a vector $\vec{b}$ makes
$\mat{A}\vec{x}=\vec{b}$ solvable iff $\vec{b}$ is orthogonal to the left nullspace. Find that
orthogonality condition explicitly and check it on $\vec{b}=\colvec{1\\2\\0}$ and
$\vec{b}=\colvec{1\\1\\1}$.
\end{exercise}

\begin{exercise}\label{ch15:ex-rowcol-iso}
Let $\mat{A}=\begin{bmatrix}2&0\\0&3\\0&0\end{bmatrix}$. Identify all four subspaces, then verify
by hand that $\mat{A}$ carries the row-space basis to the column-space basis in a one-to-one onto
fashion (Theorem~\ref{ch15:iso-thm}). What are the two stretch factors?
\end{exercise}

\begin{exercise}\label{ch15:ex-proof}
\emph{(Proof.)} Prove that for any $m\times n$ matrix $\mat{A}$, a vector in both the row space and
the nullspace must be the zero vector. Then use this to prove that the only vector in all four
subspaces simultaneously is $\vec{0}$ (after accounting for the rooms they live in).
\end{exercise}

\begin{exercise}\label{ch15:ex-transpose}
\emph{(Symmetry challenge.)} Explain why $\rank(\mat{A})=\rank(\mat{A}\T)$ guarantees that the
column space and row space always have the \emph{same} dimension, even when $m\ne n$ so they live
in rooms of different sizes. Illustrate with a $2\times 5$ matrix of rank $2$: what are the four
dimensions?
\end{exercise}

\begin{exercise}\label{ch15:ex-application}
\emph{(Modeling.)} A data matrix $\mat{A}$ is $1000\times 50$ (1000 samples, 50 features) and turns
out to have rank $r=8$. In plain language, what do the column space, nullspace, row space, and left
nullspace each mean for this dataset? Which subspace measures ``directions in feature space that do
not affect any output,'' and which measures ``output patterns the data can never produce''?
\end{exercise}
